\section{Verifica e risultati}
\label{sec:risultati}

\subsection{Metodo}
\label{sec:metodo}

Ogni caso del lavoro di riferimento è stato calcolato con il programma
zero-dimensionale corrispondente; i casi delle fig.~3a, 3b, 4, 5 e 7 sono
stati calcolati anche con il solver, su una mesh di una sola cella, con le
stesse condizioni iniziali, gli stessi dati e lo stesso passo di un
nanosecondo. I risultati sono sottoposti a tre verifiche:

\begin{enumerate}[leftmargin=1.8em]
  \item \textbf{solver contro programmi zero-dimensionali}: i due percorsi
    fanno la stessa fisica, quindi ogni differenza viene dal modo in cui il
    solver integra il modello; è la verifica dell'implementazione;
  \item \textbf{programmi zero-dimensionali contro il lavoro di riferimento}:
    è la verifica della fisica;
  \item \textbf{temperatura finale contro conservazione dell'energia}: la
    temperatura raggiunta all'equilibrio deve essere quella calcolata
    dall'energia iniziale, e quella indicata nel lavoro di riferimento.
\end{enumerate}

Le curve di \hyfoam{} sono state ricavate direttamente dal pdf dell'articolo.
Le figure sono disegni vettoriali: ogni curva è memorizzata come una sequenza
di punti, che è stata letta e riportata in tempo e temperatura usando le tacche
degli assi. L'incertezza è circa lo spessore della linea, 40--60~K sulle
temperature. Gli scarti sulle temperature sono dati in percentuale, come
massimo lungo il transitorio. Le densità delle specie sono disegnate in scala
logaritmica, e il loro scarto è dato in \emph{decadi}: una decade è un fattore
10, e 0{,}01 decadi corrispondono a circa il 2{,}3\%. I programmi salvano i
risultati ogni 10 o 100~ns, e i confronti partono da lì. Nei grafici le linee
chiare continue sono il lavoro di riferimento (\emph{paper}), quelle
tratteggiate i programmi zero-dimensionali (\emph{standalone}) e i punti il
solver (\emph{solver}).

\subsection{Azoto puro}
\label{sec:risultati-azoto}

Le fig.~3a e 3b del lavoro di riferimento riguardano l'azoto a 1~atm, riscaldato
nella vibrazione (10\,000~K per la traslazione, 1000~K per la vibrazione) o
raffreddato (3000~K e 10\,000~K). Nel riscaldamento la temperatura finale è
7623{,}3~K, esattamente il valore dato nel testo dell'articolo, e le curve
differiscono al massimo dello 0{,}3\% sulla temperatura traslazionale e
dell'1{,}5\% su quella vibrazionale. Nel raffreddamento la temperatura finale è
4973~K contro 4972~K; lo scarto nel transitorio sale al 2{,}8\% e al 5{,}1\%,
e corrisponde a un piccolo ritardo del rilassamento
(figura~\ref{fig:azoto-3}).

\begin{figure}[ht]
  \centering
  \begin{subfigure}{0.49\textwidth}
    \includegraphics[width=\linewidth]{fig3a}
    \caption{Riscaldamento vibrazionale (fig.~3a).}
  \end{subfigure}\hfill
  \begin{subfigure}{0.49\textwidth}
    \includegraphics[width=\linewidth]{fig3b}
    \caption{Raffreddamento vibrazionale (fig.~3b).}
  \end{subfigure}
  \caption{Azoto puro a 1~atm: temperatura traslazionale ($T_{tr}$) e
  vibrazionale ($T_v$) nel tempo.}
  \label{fig:azoto-3}
\end{figure}

La fig.~4 parte da 30\,000~K e mette a confronto il caso senza e con energia
elettronica (figura~\ref{fig:azoto-4}). Le temperature finali sono 21\,932~K
e 17\,650~K, contro i 21{,}9 e 17{,}7~kK indicati in figura: l'energia
elettronica abbassa la temperatura finale di oltre 4000~K. Gli scarti nel
transitorio sono dello 0{,}05\% e dello 0{,}8\% senza energia elettronica,
dello 0{,}3\% e del 4\% con. Questo caso giustifica la scelta di riscrivere lo
scambio V--T con l'energia vibro-elettronica
(sezione~\ref{sec:implementato-limiti}): con la forma originale di \mpp, che
usa la sola energia vibrazionale, la temperatura finale sarebbe la stessa ma
gli scarti nel transitorio salirebbero al 28\% e al 56\%.

\begin{figure}[ht]
  \centering
  \begin{subfigure}{0.49\textwidth}
    \includegraphics[width=\linewidth]{fig4-noEl}
    \caption{Senza energia elettronica.}
  \end{subfigure}\hfill
  \begin{subfigure}{0.49\textwidth}
    \includegraphics[width=\linewidth]{fig4-el}
    \caption{Con energia elettronica.}
  \end{subfigure}
  \caption{Azoto puro da 30\,000~K (fig.~4); la linea punteggiata orizzontale
  è la temperatura finale attesa dalla conservazione dell'energia.}
  \label{fig:azoto-4}
\end{figure}

\subsection{Miscele non reagenti}
\label{sec:risultati-miscele}

La fig.~5 riguarda una miscela di azoto molecolare e atomico, con
$5\cdot10^{22}$ particelle per metro cubo di ciascuna specie, a 30\,000~K e
1000~K. La temperatura finale è 24\,353~K contro i 24{,}4~kK della figura, e la
temperatura traslazionale resta entro il 2{,}5\%; quella vibrazionale, però,
sale più tardi che nel lavoro di riferimento, con uno scarto che arriva al
20\% nel tratto più ripido. L'origine più probabile è il tempo di
rilassamento nelle miscele, che \mpp{} media in modo diverso dal lavoro di
riferimento (sezione~\ref{sec:implementato-limiti}): il rilassamento
vibrazionale risulta un po' più lento, mentre lo stato finale è corretto.

La fig.~6 riguarda una miscela di azoto e ossigeno in parti uguali, a 1~atm,
con la vibrazione molto più calda della traslazione (30\,000~K contro 5000~K),
calcolata senza e con lo scambio V--V fra le due molecole. La temperatura
finale è 12\,133~K contro i 12{,}1~kK della figura. Senza scambio V--V gli
scarti sono del 2{,}7\% sulla temperatura traslazionale, del 2{,}7\% su quella
vibrazionale dell'azoto e del 6{,}4\% su quella dell'ossigeno; con lo scambio
scendono al 2{,}0\%, 2{,}5\% e 3{,}4\%. Come nel lavoro di riferimento, lo
scambio V--V avvicina le due temperature vibrazionali e rende più rapido il
ritorno all'equilibrio (figura~\ref{fig:miscele}).

\begin{figure}[ht]
  \centering
  \begin{subfigure}{0.49\textwidth}
    \includegraphics[width=\linewidth]{fig5}
    \caption{Azoto molecolare e atomico (fig.~5).}
  \end{subfigure}\hfill
  \begin{subfigure}{0.49\textwidth}
    \includegraphics[width=\linewidth]{fig6-VV}
    \caption{Azoto e ossigeno, con scambio V--V (fig.~6).}
  \end{subfigure}
  \caption{Miscele non reagenti. Nella miscela di azoto e ossigeno ogni
  molecola ha la propria temperatura vibrazionale.}
  \label{fig:miscele}
\end{figure}

\subsection{Miscele reagenti}
\label{sec:risultati-reagenti}

Le fig.~7 e 8 riprendono la miscela di azoto molecolare e atomico, ora con la
dissociazione dell'azoto. Nella fig.~7, che parte da 30\,000~K e 1000~K, alla
fine del calcolo le due temperature valgono 8309~K e 8312~K contro 8292~K e
8281~K del lavoro di riferimento, e le densità dell'azoto molecolare e atomico,
riferite a quella iniziale, valgono 0{,}138 e 1{,}225 contro 0{,}139 e 1{,}222:
la composizione segue quella del lavoro di riferimento entro 0{,}006 decadi,
cioè l'1{,}4\%, per tutto il calcolo (figura~\ref{fig:reagenti}). La
temperatura traslazionale resta entro il 2{,}3\%; quella vibrazionale mostra
lo stesso ritardo della fig.~5, fino al 17\% all'inizio. La fig.~8 parte
invece da 30\,000~K per entrambe le temperature; le temperature finali sono
12\,121~K e 12\,141~K contro 12\,101~K e 12\,043~K, la composizione è entro
0{,}01 decadi e le temperature entro il 2\% e il 10\%.

Questi casi confermano anche la scelta dell'esponente della temperatura di
Park. Con il valore fisso di \mpp{} ($a = 0{,}5$) la composizione della fig.~7
si allontana dal lavoro di riferimento fino a 0{,}038 decadi (9\%) e la
temperatura traslazionale fino al 6{,}7\%; con $a = 0{,}7$, come nel lavoro di
riferimento, gli scarti scendono a 0{,}006 decadi e al 2{,}3\%. La temperatura
vibrazionale non cambia, perché il suo scarto viene dal rilassamento e non
dalla chimica.

\begin{figure}[ht]
  \centering
  \begin{subfigure}{0.49\textwidth}
    \includegraphics[width=\linewidth]{fig7}
    \caption{Temperature.}
  \end{subfigure}\hfill
  \begin{subfigure}{0.49\textwidth}
    \includegraphics[width=\linewidth]{fig7-n}
    \caption{Densità delle specie riferite a quella iniziale.}
  \end{subfigure}
  \caption{Miscela reagente di azoto molecolare e atomico da 30\,000~K e
  1000~K (fig.~7), con esponente di Park $a = 0{,}7$.}
  \label{fig:reagenti}
\end{figure}

\subsection{Aria a cinque specie}
\label{sec:risultati-aria}

La fig.~9 riguarda aria (79\% di azoto e 21\% di ossigeno in moli) a
0{,}063~atm e 10\,000~K, con diciannove reazioni e una sola temperatura. Con le
costanti Quantum-Kinetic la temperatura segue il lavoro di riferimento entro lo
0{,}7\% (5096~K contro 5074~K alla fine) e le densità delle cinque specie entro
0{,}06 decadi, quattro su cinque entro 0{,}03 (figura~\ref{fig:aria}). Con il
meccanismo di Park di \mpp{}, invece, la temperatura si scosta fino all'8{,}9\%
e le densità di ossigeno molecolare e atomico, azoto atomico e monossido di
azoto di 0{,}2--0{,}5 decadi, cioè di un fattore fra 1{,}6 e 3: la fig.~9 del
lavoro di riferimento è stata ottenuta con le costanti Quantum-Kinetic.

\begin{figure}[ht]
  \centering
  \begin{subfigure}{0.49\textwidth}
    \includegraphics[width=\linewidth]{fig9-QK}
    \caption{Temperatura.}
  \end{subfigure}\hfill
  \begin{subfigure}{0.49\textwidth}
    \includegraphics[width=\linewidth]{fig9-QK-n}
    \caption{Densità delle specie riferite a quella iniziale.}
  \end{subfigure}
  \caption{Aria a cinque specie da 10\,000~K con le costanti Quantum-Kinetic
  (fig.~9).}
  \label{fig:aria}
\end{figure}

\subsection{Solver e programmi zero-dimensionali}
\label{sec:risultati-solver}

La tabella~\ref{tab:solver} riporta lo scarto massimo fra il solver e i
programmi zero-dimensionali in tutti i casi calcolati con entrambi. Le
temperature finali coincidono, e le curve restano entro lo 0{,}2\% in tutti i
casi tranne la fig.~4, dove lo scarto arriva all'1{,}2\% nel rilassamento
molto rapido dei primi 100~ns. Lì conta di più il modo in cui si avanza nel
tempo: il solver ripete il calcolo due volte a ogni passo, i programmi
zero-dimensionali una sola. Anche con le reazioni accese la composizione
coincide entro lo 0{,}03\%. Il modello è quindi integrato correttamente nel
solver, chimica compresa.

\begin{table}[ht]
  \centering
  \caption{Scarto massimo fra solver e programmi zero-dimensionali.}
  \label{tab:solver}
  \small
  \begin{tabular}{lccc}
    \toprule
    Caso & $T_{tr}$ & $T_{v}$ & Composizione \\
    \midrule
    fig.~3b, raffreddamento            & 0{,}00\% & 0{,}00\% & -- \\
    fig.~3a, riscaldamento             & 0{,}03\% & 0{,}21\% & -- \\
    fig.~4, senza energia elettronica  & 0{,}14\% & 0{,}69\% & -- \\
    fig.~4, con energia elettronica    & 0{,}51\% & 1{,}18\% & -- \\
    fig.~5, azoto molecolare e atomico & 0{,}02\% & 0{,}13\% & -- \\
    fig.~7, con reazioni               & 0{,}03\% & 0{,}18\% & 0{,}03\% \\
    \bottomrule
  \end{tabular}
\end{table}
